\documentclass{situationist-preprint}

\title{Fair Play and Political Obligation}
\msid{POL-03-b}
\msfamily{POL-03}
\msdiscipline{Political Philosophy}
\mstype{Companion manuscript}
\mskeywords{political obligation; fair play; particularity; discretionary goods; Hart; Klosko; Simmons}

\begin{document}
\maketitle

\begin{abstract}
POL-03-a argues that there is a weak, defeasible duty to obey just laws, grounded in fairness. We test the fair-play account against two further requirements that Simmons sets for any theory of political obligation and that POL-03-a does not discuss: particularity, which asks why the duty is owed to one's own state rather than to just states generally, and the problem of discretionary goods, which asks how a principle grounded in presumptively beneficial goods can bind us to laws that provide goods some do not want. We argue that fair play meets the first more naturally than its rivals, and that it meets the second once the duty is located in the scheme of just law as a whole, at the strength POL-03-a assigns it.
\end{abstract}

\section{Two requirements}

\citet{Simmons1979} argues that an account of political obligation should satisfy a requirement of \emph{particularity}: it should explain why a person is bound to the laws of her own state in a way in which she is not bound to the laws of other states, however just. And critics of the principle of fairness since \citet{Nozick1974} have pressed the problem of \emph{discretionary goods}: a principle grounded in receipt of benefits can bind only to the extent that the benefits are wanted, and many of the goods that laws provide, from public broadcasting to the arts, are goods on which reasonable persons differ \citep{Klosko1992}.

The account of POL-03-a rests on the principle of fairness first stated by \citet{Hart1955} (Principle~P1 of POL-03-a), the assumption that the legal order of a reasonably just state is a cooperative scheme whose goods are presumptively beneficial (Assumption~2.3), and the marginality of single breaches (Assumption~2.4). From these it derives a duty that exists (Lemma~3.1), is weak (Lemma~3.2) and is defeasible (Lemma~3.3).

\section{Particularity}

Natural-duty accounts derive the duty to obey from a general duty to support just institutions \citep[\S 19]{Rawls1971}. They have difficulty with particularity: the duty to support just institutions is owed to just institutions wherever they are, and a further premise is needed to explain why my own state has a special claim on me.

\begin{proposition}\label{p:particular}
The duty grounded in fairness is owed to the participants in the scheme whose goods one receives, and so, in the ordinary case, to one's fellow subjects of one's own state.
\end{proposition}

\begin{proof}
By Principle~P1 of POL-03-a, the duty is owed by a person who receives goods produced by a scheme to the other participants in that scheme. The goods of security, the rule of law and coordination are produced by the compliance of those subject to a particular legal order and received by those within its jurisdiction. A person receives them from her own state's scheme and not from the schemes of other states.
\end{proof}

Particularity thus falls out of the structure of the principle rather than requiring an additional premise. It also explains the duties of resident aliens and visitors, who receive the goods of the scheme while within it, and their absence once they leave, which is what ordinary judgement holds.

\section{Discretionary goods}

The second problem is harder. If the duty rests on goods being worth their cost to everyone, it seems to bind only to those laws that provide such goods, and to leave a person free to ignore laws that provide goods she does not value.

We reply that the duty attaches to the scheme of just law as a whole, of which the laws providing discretionary goods are part. A reasonably just legal order is not a set of independent schemes, one per law, but a single scheme in which the presumptively beneficial goods are produced by general compliance with a system of rules that also includes rules on contested matters, adopted by fair procedures. The person who receives the security and coordination of the system receives them from a scheme whose terms include that contested matters are settled by those procedures. Her share of the burden is compliance with the system's just laws, including those whose particular benefits she does not value.

This reply has a cost that must be acknowledged: the duty with respect to discretionary laws is no stronger than the duty with respect to the scheme, and that duty is weak. By Lemma~3.2 of POL-03-a, the wrong of a single breach is the wrong of free-riding on a fair scheme, and reasons of moderate strength outweigh it. This is the right result. A person's conscientious objection to a law providing a good she rejects is a reason of the relevant strength, and the account permits it to prevail in many cases, while denying that indifference to the good is enough.

\section{Objections}

\begin{objection}
Treating the legal order as one scheme lets the presumptively beneficial goods launder any law, however objectionable.
\end{objection}

\begin{reply}
It does not, because the duty is to obey \emph{just} laws of a \emph{reasonably just} scheme (Definition~2.2 of POL-03-a). For seriously unjust laws the duty lapses (Lemma~3.3 of POL-03-a), since compliance with them is not a fair share of a fair burden. Under an order that is unjust in some respects but not extremely so, much of the law remains binding, and the laws that are seriously unjust do not.
\end{reply}

\begin{objection}
Particularity is secured at the price of generality: those who receive little from the scheme, such as the marginalized, owe little or nothing.
\end{objection}

\begin{reply}
The goods of security and the rule of law are received by nearly everyone subject to a reasonably just order, and the account does imply that those to whom the scheme denies them owe less. Where whole groups are so denied, the scheme is not reasonably just, and the conditions of the duty fail for them. This is a consequence the account should have.
\end{reply}

\section{Conclusion}

The fair-play account explains why the duty to obey is owed to one's own state, and it extends to laws that provide discretionary goods by locating the duty in the scheme of just law as a whole. In both respects the duty remains what POL-03 concludes it to be: weak, defeasible, owed only with respect to just laws, and grounded in fairness.

\begin{thebibliography}{9}
\bibitem[Hart(1955)]{Hart1955} Hart, H.~L.~A. (1955). Are there any natural rights? \emph{The Philosophical Review}, 64(2), 175--191.
\bibitem[Klosko(1992)]{Klosko1992} Klosko, G. (1992). \emph{The Principle of Fairness and Political Obligation}. Rowman \& Littlefield.
\bibitem[Nozick(1974)]{Nozick1974} Nozick, R. (1974). \emph{Anarchy, State, and Utopia}. Basic Books.
\bibitem[Rawls(1971)]{Rawls1971} Rawls, J. (1971). \emph{A Theory of Justice}. Harvard University Press.
\bibitem[Simmons(1979)]{Simmons1979} Simmons, A.~J. (1979). \emph{Moral Principles and Political Obligations}. Princeton University Press.
\end{thebibliography}

\end{document}
